% ============================================================================
% CAPÍTULO 5: ESTUDIO DE FACTIBILIDAD ECONÓMICA
% ============================================================================

\chapter{Estudio de Factibilidad Económica}
\label{cap:factibilidad_economica}

Para garantizar la viabilidad y sostenibilidad a largo plazo del sistema de control de acceso vehicular automatizado propuesto para el Campus Ñuñoa de la Universidad Tecnológica Metropolitana, es imperativo complementar el diseño de la arquitectura técnica con un análisis financiero riguroso. Este capítulo desarrolla la evaluación económica del proyecto estructurada bajo la metodología de modernización o \textit{Retrofit}.

El enfoque de \textit{Retrofit} consiste en reutilizar y complementar la infraestructura preexistente de la institución (como el portón metálico corredero, la cremallera y los sistemas de motorización eléctrica existentes) y superponer una capa inteligente de hardware de borde y software de visión computacional. Este paradigma permite mitigar drásticamente los costos de inversión iniciales al prescindir de costosas obras de edificación civil o renovación de motores.

La evaluación se proyecta sobre un horizonte temporal de \textbf{5 años}. De acuerdo con la Metodología General del Ministerio de Desarrollo Social y Familia \parencite{METODOLOGIAGENERAL2025}, el horizonte de evaluación depende principalmente de la vida útil de los bienes de capital principales y de la previsibilidad del proyecto. Dado que el sistema se fundamenta en hardware de Inteligencia Artificial sujeto a una rápida obsolescencia y dada la baja previsibilidad de estas tecnologías a largo plazo, se adopta este plazo \textit{ad hoc} como el más conservador y adecuado. Esta decisión evita modelar complejas reinversiones, capturando el remanente de vida útil contable de los equipos (6 y 9 años según el SII) a través de la recuperación de un valor residual al término del quinto año. Para determinar la rentabilidad institucional, se detallan los costos de inversión de capital (CAPEX), costos operativos anuales (OPEX), la depreciación según normativa del Servicio de Impuestos Internos (SII), la recuperación de valores residuales y los indicadores financieros de rentabilidad (Valor Actual Neto, VAN; Tasa Interna de Retorno, TIR; y Periodo de Recuperación de Capital, PRC) calculados bajo la Tasa Social de Descuento (TSD) del \textbf{5,5\%} anual fijada por el MIDESO \parencite{MinisterioDesarrolloSocial}.

\section{Inversión Inicial (CAPEX)}
\label{sec:capex}

La inversión inicial o costos de capital (CAPEX, \textit{Capital Expenditure}) contempla la adquisición de todos los componentes de hardware, los materiales físicos de canalización y montaje, y la mano de obra especializada requerida para la instalación y puesta en marcha del sistema en el Campus Ñuñoa. Al tratarse de una solución institucional del sector público, todos los precios se consignan en valores brutos con IVA (19\%) incluido, dado que la universidad actúa como consumidor final y no recupera crédito fiscal.

A continuación, en la Tabla~\ref{tab:capex}, se detalla el presupuesto estructurado del CAPEX para el proyecto de modernización tecnológica, agrupado en componentes lógicos, infraestructura pasiva y costos de implementación.

\begin{table}[H]
    \centering
    \small
    \setlength{\tabcolsep}{4.5pt}
    \caption{Presupuesto de Inversión Inicial (CAPEX) para el sistema de control de acceso.}
    \begin{tabular}{|p{5.2cm}|p{3.2cm}|c|r|}
        \hline
        \textbf{Componente / Item} & \textbf{Proveedor} & \textbf{Cantidad} & \textbf{Costo Total (CLP)} \\ \hline
        \multicolumn{4}{|l|}{\textbf{1. Equipos de Cómputo y Visión (Capa Inteligente)}} \\ \hline
        NVIDIA Jetson Orin Nano 8GB & MCI Electronics & 1 & \$800.000 \\ \hline
        Cámara IP LPR Dahua ITC413 & TiCam & 2 & \$1.476.790 \\ \hline
        Módulo Relé Optoacoplado 5V & MechatronicStore & 1 & \$2.300 \\ \hline
        Gabinete Metálico IP65 250x200x150 & Lexo Chile & 1 & \$31.500 \\ \hline
        Switch PoE Gigabit Dahua & SmartSecure & 1 & \$55.835 \\ \hline
        \multicolumn{4}{|l|}{\textbf{2. Infraestructura Pasiva y Montaje}} \\ \hline
        Pedestal Metálico CCTV 3m Cónico & LK Distribuidor & 1 & \$147.084 \\ \hline
        Cable UTP Cat6 Exterior (30m) & TiCam & 1 & \$23.220 \\ \hline
        Cable SPT Paralelo 2x18 AWG (20m) & Sodimac & 1 & \$27.800 \\ \hline
        Tuberías y Accesorios PVC Conduit & Sodimac & 1 (Lote) & \$8.308 \\ \hline
        Hormigón H20 Topex (Saco 25kg) & Sodimac & 5 & \$14.300 \\ \hline
        \multicolumn{4}{|l|}{\textbf{3. Implementación y Puesta en Marcha}} \\ \hline
        Mano de Obra Especializada & Servicio Externo & 1 & \$480.000 \\ \hline
        \hline
        \multicolumn{3}{|l|}{\textbf{CAPEX TOTAL}} & \textbf{\$3.067.137} \\ \hline
    \end{tabular}
    \label{tab:capex}
\end{table}

Este CAPEX total de \textbf{\$3.067.137 CLP} se compone en un 84,35\% por la adquisición directa de activos tangibles e infraestructura de hardware (\$2.587.137 CLP) respaldada en cotizaciones de mercado (véase Anexo~\ref{anexo:cotizaciones}), y en un 15,65\% por los servicios profesionales de montaje físico y cableado (\$480.000 CLP).

El cálculo de la \textbf{Mano de Obra Especializada} se fundamenta en las tarifas de mercado para técnicos instaladores de corrientes débiles y sistemas CCTV en la Región Metropolitana. Se determinó la necesidad de una cuadrilla compuesta por 2 técnicos trabajando de manera simultánea durante 3 jornadas completas para el montaje de la obra civil menor, canalización y configuración del sistema de borde. Según ofertas laborales vigentes (Jooble Chile, 2026; véase Anexo~\ref{anexo:cotizaciones}), el costo referencial asciende a \$80.000 CLP diarios por técnico, resultando en un costo total de instalación de \$480.000 CLP (2 técnicos $\times$ 3 días $\times$ \$80.000 CLP/día).

\section{Costos Operativos Anuales (OPEX)}
\label{sec:opex}

Los gastos de operación (OPEX, \textit{Operational Expenditure}) representan los costos recurrentes que la institución debe absorber anualmente para mantener el sistema operativo de forma ininterrumpida y segura. Para esta solución de bajo consumo eléctrico y alta robustez física, el OPEX anual se compone únicamente de dos variables críticas:

\begin{enumerate}
    \item \textbf{Consumo Eléctrico de Borde:} Corresponde a la energía eléctrica demandada por la NVIDIA Jetson Orin Nano, el Switch PoE y las dos cámaras IP LPR. Basado en las fichas técnicas de los fabricantes (véase Anexo~\ref{anexo:cotizaciones}), se adopta un modelo de carga mixta: 12 horas diurnas en estado nominal ($\sim$25 W) y 12 horas nocturnas a máxima potencia con iluminación activa ($\sim$62 W), resultando en un consumo diario de 1,044 kWh. Conforme a la Resolución Exenta N° 03007 de la \textcite{UTEMRes30072021} y el Pliego Tarifario Regulado de \textcite{EnelDistribucionChile} para la comuna de Ñuñoa, se aplica la opción tarifaria institucional AT3 (Media/Alta Tensión), con un costo de energía actualizado de \textbf{\$147,63 CLP por kWh}. El gasto anual se calcula como:
    \begin{align}
        \text{Consumo Anual} = & \, 1{,}044 \text{ kWh/día} \times 365 \text{ días/año} \nonumber \\
        & \times \$147{,}63 \text{ CLP/kWh} = \$56.256 \text{ CLP}
    \end{align}
    Para efectos del análisis presupuestario, este valor se redondea a \textbf{\$56.260 CLP} anuales.
    
    \item \textbf{Mantenimiento Preventivo y Limpieza:} Contempla dos visitas técnicas semestrales especializadas para efectuar la limpieza de lentes ópticos (eliminando acumulación de hollín y polvo ambiental), reapriete de borneras eléctricas en gabinete, aplicación de limpia-contactos dieléctrico e inspección de estanqueidad IP65. Se modela en base a la misma tarifa de mercado de técnico instalador CCTV (\$10.000 CLP/hh; Jooble Chile, 2026; véase Anexo~\ref{anexo:cotizaciones}), requiriendo 4 horas de inspección técnica por visita (\$40.000 CLP c/u). Se proyecta un costo anual de \textbf{\$80.000 CLP} (\$40.000 CLP $\times$ 2 visitas).
\end{enumerate}

La Tabla~\ref{tab:opex} resume la estructura de costos operativos anuales proyectados para el sistema.

\begin{table}[H]
    \centering
    \caption{Costos Operativos Anuales (OPEX) proyectados.}
    \begin{tabular}{|p{8cm}|r|}
        \hline
        \textbf{Concepto de Gasto} & \textbf{Costo Anual (CLP)} \\ \hline
        Consumo Eléctrico del Nodo de Borde & \$56.260 \\ \hline
        Mantenimiento Preventivo y Limpieza Semestral & \$80.000 \\ \hline
        \hline
        \textbf{OPEX ANUAL TOTAL} & \textbf{\$136.260} \\ \hline
    \end{tabular}
    \label{tab:opex}
\end{table}

\section{Cuantificación de Beneficios y Ahorros del Proyecto}
\label{sec:ahorros}

El beneficio económico de la automatización del control de acceso no proviene de la venta de servicios ni de cobros directos a los usuarios, sino de la optimización del recurso humano disponible en la portería del campus.

En la situación actual (sin automatización), los funcionarios de conserjería deben interrumpir sus labores de seguridad, vigilancia y atención a los usuarios para verificar de manera visual a cada conductor y accionar manualmente el portón vehicular a través de control remoto. Se estima empíricamente que el tiempo dedicado a esta labor repetitiva y de bajo valor agregado suma un total de \textbf{2 horas diarias} por jornada de operación del campus.

La implementación del sistema ALPR automatiza por completo este proceso de verificación, eliminando la necesidad de intervención manual para el flujo de vehículos de cuatro ruedas empadronados dentro del alcance definido del proyecto (véase la \hyperref[sec:alcances_limitaciones]{sección de Alcances y Limitaciones}). De esta forma, el conserje puede reasignar estas 2 horas diarias a tareas críticas de seguridad activa y control de perímetros.

Para cuantificar económicamente este beneficio institucional, se modela la reasignación de tiempo bajo los siguientes supuestos:
\begin{itemize}
    \item \textbf{Jornada Mensual Operativa:} El campus opera bajo flujos regulares de acceso durante 22 días hábiles al mes.
    \item \textbf{Horizonte Anual:} 12 meses de funcionamiento administrativo continuo.
    \item \textbf{Valorización de Hora-Hombre (HH):}\footnote{Valorización calculada en base al Ingreso Mínimo Mensual legal vigente (Ley N°21.830 \parencite{LeySueldoMinimo2026}), incorporando los costos de empresa obligatorios (leyes sociales y seguros) para una jornada estandarizada de 180 horas mensuales. El desglose matemático exacto se detalla en el Anexo~\ref{anexo:cotizaciones}.} Se estima un costo hora de \textbf{\$4.000 CLP} para el personal de portería.
\end{itemize}

El ahorro financiero anual por concepto de optimización del recurso humano se calcula de la siguiente manera:
\begin{align}
    \text{Ahorro Anual} = & \, 2 \text{ h/día} \times 22 \text{ días/mes} \times 12 \text{ meses/año} \nonumber \\
    & \times \$4.000 \text{ CLP/HH} = \$2.112.000 \text{ CLP}
\end{align}

Este monto representa un flujo de beneficio indirecto o ahorro operativo de \textbf{\$2.112.000 CLP} anuales, el cual ingresa al flujo de caja del proyecto como el principal retorno financiero de la inversión inicial.

\section{Depreciación Contable y Valor Residual de Activos (Norma SII)}
\label{sec:depreciacion_sii}

Para evitar arbitrariedades en la determinación de la vida útil de los activos y respaldar rigurosamente el análisis económico, la depreciación contable se fundamenta en las \textbf{Tablas de Vida Útil de los Bienes Físicos del Activo Fijo} fijadas por el \textbf{Servicio de Impuestos Internos (SII)} de Chile \parencite[Resolución Exenta N° 43 del 2002]{SIIServicioImpuestos} y las pautas metodológicas de formulación de proyectos MIDESO:

\begin{enumerate}
    \item \textbf{Equipos de Procesamiento y Cómputo (Jetson Orin Nano, Switch PoE y Relé):} Clasificados según la normativa del SII en la categoría de \textit{``Sistemas de procesamiento de datos, computadores y redes''}, a los cuales se les asigna una vida útil normal de \textbf{6 años}. Subtotal: \$858.135 CLP. Depreciación anual: \$143.023 CLP/año.
    \item \textbf{Óptica y Seguridad CCTV (Cámaras LPR Dahua ITC413):} Clasificadas en la categoría de \textit{``Equipos de televisión en circuito cerrado (CCTV) e instrumentación electrónica''}. Dada su naturaleza de equipos electrónicos de intemperie para operación industrial continua (protección IP67 e IK10), se adopta la vida útil normal del SII de \textbf{9 años}. Subtotal: \$1.476.790 CLP. Depreciación anual: \$164.088 CLP/año.
    \item \textbf{Infraestructura Pasiva y Obras Civiles (Pedestal, Gabinete IP65, Cableado UTP/Control, Conduit PVC y Hormigón H20):} Clasificados en la categoría de \textit{``Construcciones, obras civiles y estructuras metálicas livianas''}, con una vida útil normal de \textbf{10 años}. Subtotal: \$252.212 CLP. Depreciación anual: \$25.221 CLP/año.
\end{enumerate}

Como todas las vidas útiles de los componentes superan el horizonte de evaluación de 5 años del proyecto, \textbf{no se requiere modelar una reinversión de equipos}.

Al término del período de evaluación (Año 5), los activos no han concluido su vida útil física ni contable, por lo que conservan un \textbf{Valor Libro Remanente (Valor Residual)}. Este valor se calcula restando la depreciación acumulada de 5 años al costo original de los equipos y se rescata como un ingreso de capitales en el flujo del Año 5.

La Tabla~\ref{tab:depreciacion} resume la amortización contable y el cálculo exacto del Valor Residual del proyecto.

\begin{table}[H]
    \centering
    \footnotesize
    \setlength{\tabcolsep}{3.5pt}
    \caption{Depreciación contable y Valor Residual de activos al Año 5 (Normativa SII).}
    \begin{tabular}{|p{4.2cm}|c|r|r|r|}
        \hline
        \textbf{Categoría del Activo (SII)} & \textbf{Vida Útil} & \textbf{Costo Bruto} & \textbf{Deprec. Anual} & \textbf{Valor Residual (Año 5)} \\ \hline
        Sistemas de Cómputo (Jetson/Switch) & 6 años & \$858.135 & \$143.023 & \$143.020 \\ \hline
        Óptica y CCTV (Cámaras Dahua) & 9 años & \$1.476.790 & \$164.088 & \$656.350 \\ \hline
        Infraestructura Pasiva y Estructura & 10 años & \$252.212 & \$25.221 & \$126.110 \\ \hline
        \hline
        \textbf{TOTALES} & -- & \textbf{\$2.587.137} & \textbf{\$332.332} & \textbf{\$925.480} \\ \hline
    \end{tabular}
    \label{tab:depreciacion}
\end{table}

\section{Flujo de Caja Proyectado y Evaluación Económica}
\label{sec:flujo_caja}

Siguiendo la estructura formal de evaluación de proyectos, el flujo financiero se desglosa en el \textbf{Flujo de Caja Operacional} (beneficios operativos menos OPEX) y el \textbf{Flujo de Caja de Capitales} (inversión inicial CAPEX en Año 0 y la recuperación del Valor Residual en el Año 5). 

En el contexto de una universidad pública, la evaluación se realiza exenta de Impuesto a la Renta de Primera Categoría (al no constituir una actividad comercial gravada), por lo que la depreciación no genera escudo fiscal pero permite determinar con precisión la recuperación del activo fijo.

La Tabla~\ref{tab:flujo_caja} muestra la estructura del Flujo de Caja Neto Proyectado a 5 años.

\begin{table}[H]
    \centering
    \footnotesize
    \setlength{\tabcolsep}{2pt}
    \caption{Estructura del Flujo de Caja Neto Proyectado a 5 años (CLP).}
    \begin{tabular}{|p{2.8cm}|r|r|r|r|r|r|}
        \hline
        \textbf{Concepto / Año} & \textbf{Año 0} & \textbf{Año 1} & \textbf{Año 2} & \textbf{Año 3} & \textbf{Año 4} & \textbf{Año 5} \\ \hline
        \multicolumn{7}{|l|}{\textbf{1. Flujo de Caja Operacional (FCO)}} \\ \hline
        Ahorro Operativo (Conserjería) & \$0 & \$2.112.000 & \$2.112.000 & \$2.112.000 & \$2.112.000 & \$2.112.000 \\ \hline
        Costos de Operación (OPEX) & \$0 & -\$136.260 & -\$136.260 & -\$136.260 & -\$136.260 & -\$136.260 \\ \hline
        \textbf{Subtotal Flujo Operacional} & \textbf{\$0} & \textbf{\$1.975.740} & \textbf{\$1.975.740} & \textbf{\$1.975.740} & \textbf{\$1.975.740} & \textbf{\$1.975.740} \\ \hline
        \multicolumn{7}{|l|}{\textbf{2. Flujo de Caja de Capitales (FKK)}} \\ \hline
        Inversión Inicial (CAPEX) & -\$3.067.137 & \$0 & \$0 & \$0 & \$0 & \$0 \\ \hline
        Recuperación Valor Residual & \$0 & \$0 & \$0 & \$0 & \$0 & \$925.480 \\ \hline
        \textbf{Subtotal Flujo Capitales} & \textbf{-\$3.067.137} & \textbf{\$0} & \textbf{\$0} & \textbf{\$0} & \textbf{\$0} & \textbf{\$925.480} \\ \hline
        \hline
        \textbf{FLUJO DE CAJA NETO} & \textbf{-\$3.067.137} & \textbf{\$1.975.740} & \textbf{\$1.975.740} & \textbf{\$1.975.740} & \textbf{\$1.975.740} & \textbf{\$2.901.220} \\ \hline
        \textbf{FLUJO DE CAJA ACUMULADO} & \textbf{-\$3.067.137} & \textbf{-\$1.091.397} & \textbf{\$884.343} & \textbf{\$2.860.083} & \textbf{\$4.835.823} & \textbf{\$7.737.043} \\ \hline
    \end{tabular}
    \label{tab:flujo_caja}
\end{table}

\subsection{Indicadores de Rentabilidad Financiera (VAN, TIR y PRC)}

Para validar la rentabilidad y fundamentar la toma de decisiones, se calculan el Valor Actual Neto (VAN), la Tasa Interna de Retorno (TIR) y el Periodo de Recuperación de Capital (PRC) sobre el Flujo de Caja Neto, considerando la Tasa Social de Descuento del $r = 5{,}5\%$ anual vigente del Ministerio de Desarrollo Social y Familia \cite{MinisterioDesarrolloSocial}.

\begin{enumerate}
    \item \textbf{Valor Actual Neto (VAN):} Expresa la riqueza neta actualizada generada por el proyecto:
    \begin{align}
        \text{VAN} = & \, -\$3.067.137 + \sum_{t=1}^{4} \frac{\$1.975.740}{(1{,}055)^t} + \frac{\$2.901.220}{(1{,}055)^5} \nonumber \\
        = & \, \$6.077.951 \text{ CLP}
    \end{align}
    Puesto que $\text{VAN} = \$6.077.951 \text{ CLP} > 0$, el proyecto es altamente rentable e incrementa el patrimonio institucional en más de 6 millones de pesos.
    
    \item \textbf{Tasa Interna de Retorno (TIR):} Representa el rendimiento promedio anual del capital invertido:
    \begin{equation}
        0 = -\$3.067.137 + \sum_{t=1}^{4} \frac{\$1.975.740}{(1+\text{TIR})^t} + \frac{\$2.901.220}{(1+\text{TIR})^5} \implies \text{TIR} = 60{,}00\%
    \end{equation}
    Al ser la $\text{TIR} = 60{,}00\% \gg 5{,}5\%$, la inversión supera holgadamente la tasa social de exigencia del Estado.
    
    \item \textbf{Periodo de Recuperación de Capital (PRC):} Cuantifica el tiempo exacto requerido para amortizar la inversión inicial mediante la fórmula de interpolación sobre el Flujo de Caja Acumulado. Teóricamente se define como:
    \begin{equation}
        \text{PRC} = A + \frac{|FCA_{A}|}{FCN_{A+1}}
    \end{equation}
    Donde $A$ es el último período con Flujo de Caja Acumulado ($FCA$) negativo, y $FCN_{A+1}$ corresponde al Flujo de Caja Neto del período inmediatamente posterior. Sustituyendo con los valores proyectados:
    \begin{equation}
        \text{PRC} = 1 \text{ año} + \frac{|\text{-\$1.091.397}| \text{ CLP}}{\$1.975.740 \text{ CLP}} \approx 1{,}55 \text{ años}
    \end{equation}
    La inversión se recupera íntegramente en \textbf{1 año y 6 y medio meses} de funcionamiento del sistema.
\end{enumerate}

La Tabla~\ref{tab:indicadores} resume los indicadores de rentabilidad del estudio económico.

\begin{table}[H]
    \centering
    \small
    \setlength{\tabcolsep}{5pt}
    \caption{Indicadores de rentabilidad financiera del proyecto.}
    \begin{tabular}{|p{5.5cm}|r|p{3.5cm}|}
        \hline
        \textbf{Indicador Financiero} & \textbf{Valor} & \textbf{Criterio de Aceptación} \\ \hline
        Tasa Social de Descuento ($r$) & 5,50\% & Estándar Oficial MIDESO \cite{MinisterioDesarrolloSocial} \\ \hline
        Valor Actual Neto (VAN) & \$6.077.951 CLP & Aceptable (VAN > 0) \\ \hline
        Tasa Interna de Retorno (TIR) & 60,00\% & Excelente (TIR > $r$) \\ \hline
        Periodo de Recuperación (PRC) & 1,55 años & Retorno Rápido (< 3 años) \\ \hline
    \end{tabular}
    \label{tab:indicadores}
\end{table}

En conclusión, el estudio de factibilidad económica demuestra de manera objetiva y técnicamente justificada que la implementación del nodo de borde inteligente bajo la metodología \textit{Retrofit} es una inversión altamente rentable para la Universidad Tecnológica Metropolitana, garantizando el retorno del capital en el segundo año y generando un valor económico neto de más de 6 millones de pesos bajo la normativa del Sistema Nacional de Inversiones de Chile.
